\section{Historical Tanaka screening and its replacement}
\label{app:historical-cleaning}

This appendix records a preprocessing rule used for the exploratory April FNO
dataset.  The final v9 dataset was generated independently and did not inherit
this rule.  We include the details because they explain exactly which pilot
rows were removed and how that decision led to the periodic generator revision
described in \cref{sec:data}.  That replacement constructor and dataset are
complete, but the replacement dataset has not been used for reported training.

\subsection{The rule that was applied}

For one stored row, let $q^\star_j$ denote its stored DNO target at grid node
$x_j$, set $q^\star_N=q^\star_0$, and form the cyclic forward differences
$d_j=q^\star_{j+1}-q^\star_j$ for $j=0,\ldots,N-1$.  Given a relative dead
zone $\rho$, retain only the signs of differences satisfying
\begin{equation}
  |d_j|>\rho\max_\ell|d_\ell|
\end{equation}
in their cyclic spatial order.  If the remaining signs are
$s_1,\ldots,s_m$ with \(m\geq1\), define
\begin{equation}
  C_\rho(q^\star)=\sum_{i=1}^{m}\mathbf 1\{s_i\ne s_{i+1}\},
  \qquad s_{m+1}=s_1.
  \label{eq:historical-sign-count}
\end{equation}
If no difference survives the dead zone, set \(C_\rho(q^\star)=0\).
The cleaner used $\rho=0.03$ and rejected a row when
$C_{0.03}(q^\star)>10$.  It removed \num{548360} of \num{4261232} rows
($12.87\%$), leaving a \num{3712872}-row archive.  The 500,000-row Tanaka set
used by the April FNO runs was sampled from this retained archive.

\subsection{Why it detected the old generator defect}

The rule has a restricted interpretation for the pilot family.  For the
exact water-wave equations, let
$q=\G(\eta;h)\xi=\eta_t$.  An exactly translating profile
$\eta(x,t)=\bar\eta(x-ct)$ satisfies
\begin{equation}
  q=-c\bar\eta_x,
  \qquad q_x=-c\bar\eta_{xx}.
  \label{eq:traveling-curvature-count}
\end{equation}
When the stored approximation $q^\star$ resolves this physical $q$, the
dead-zoned count $C_\rho(q^\star)$ therefore approximately counts transitions
between the principal curvature lobes of a translating profile.  A smooth
unimodal crest has two inflection points.  The one-to-three-crest pilot
construction consequently has a nominal count of roughly two to six; the
cutoff at ten allowed distortion and crest interaction.

The old builder translated a real-line solitary-wave profile and interpolated
it with zero values outside the tabulated support.  A crest near a periodic
endpoint could therefore have a non-negligible tail cut to zero.  The resulting
endpoint mismatch has slowly decaying Fourier coefficients, and derivatives in
the order-six DNO amplify the associated ringing.  This produces additional
alternating curvature lobes in $q^\star$, which explains why
\cref{eq:historical-sign-count} responded to the defect.

The connection can be measured on all 254 reconstructable initial conditions
from historical batch zero.  Define the normalized discrete endpoint mismatch
\begin{equation}
  J_\eta=
  \frac{|\eta(x_0)-\eta(x_{N-1})|}{\norm{\eta}_\infty}.
  \label{eq:endpoint-mismatch}
\end{equation}
Here \(\norm{\eta}_\infty=\max_j|\eta(x_j)|\).
All 12 initial rows rejected by the sign rule had $J_\eta>0.05$, with median
$0.219$, while the 242 retained rows had median $2.50\times10^{-6}$.  For each
case, the fraction rejected over its 200 stored times had Spearman correlation
$-0.701$ with the initial distance of the nearest crest from the boundary
($p=8.4\times10^{-39}$).

A construction-level counterfactual isolates the cause.  The diagnostic
script rebuilt the same 12 rejected crest specifications by summing the
central profile and its two adjacent periodic images.  Those three images
were sufficient for these particular compactly supported profiles because
all other images vanished on $[0,L)$.  This change from nonperiodic zero
extension reduced every sign count from $12$--$38$ to $2$--$8$.  The
replacement constructor does not hard-code three images; it sums every
periodic image whose support intersects $[0,L)$.  On the three-grid fixed-band
audit of
\cref{eq:three-grid-discrepancy}, their median $N$--$2N$ DNO discrepancy falls
from $0.0827$ to $7.03\times10^{-6}$, and their median $2N$--$4N$ discrepancy
falls from $0.147$ to $2.18\times10^{-6}$.  Thus the sign rule was responding
to the old embedding operation, not merely to globally rough profiles.

\subsection{Why it is not a production validity criterion}

The count is neither necessary nor sufficient for fixed-band resolution.  At
the $10^{-3}$ two-grid tolerance, all 12 historical frame-zero sign
rejections fail the fixed-band resolution test, but so do 29 of the 242
sign-retained cases.  Thus every sign-rule rejection failed that test, but the
rule found only $12$ of its $41$ failures.  Conversely, independently
converged periodic multi-crest and Stokes states can have more than ten
physical sign transitions.
For example, three corrected periodic Tanaka states with 12, 18, and 20
transitions have fixed-band $N$--$2N$ discrepancies between
$4.90\times10^{-6}$ and $7.17\times10^{-6}$.

Consequently, the sign count is reported only as a family-specific historical
curvature screen.  The replacement construction removes its failure mechanism
by periodic wrapping and spectral restriction, while the cross-grid
fixed-band quantity \cref{eq:three-grid-discrepancy} tests the numerical
property of interest without imposing a bound on the number of physical
extrema.
